\documentclass{book}
\usepackage{hyperref}
\usepackage{amsmath}
\usepackage{CJKutf8}
\usepackage{ulem}
\usepackage{tipa} % Add this line to support \textipa command
\hypersetup{
    colorlinks=true,
    linkcolor=black,
    filecolor=magenta,      
    urlcolor=cyan,
    }

\urlstyle{same}

\title{French Pronounciation}
\author{Junzhang Luo}

\begin{document}

\tableofcontents

\chapter{Accents}

    \section{Accent Aigu (é)}
        \begin{itemize}
            \item \textipa{[e]}: \textit{été}
            \item \textipa{[je]}: \textit{épée}
        \end{itemize}

    \section{Accent Grave (è)}
        \begin{enumerate}
            \item à, ù: pronounciation doesn't change, purpose is to differentiate words
            \item è: 
                \begin{itemize}
                    \item \textipa{[E]}: \textit{père}
                \end{itemize}
        \end{enumerate}
    
    \section{Accent Circonflexe (ê)}
        \begin{enumerate}
            \item ê, same pronounciation as è
            \item â, longer and softer compare to a, \textipa{[a]}: \textit{pâte}. In real life, word with â takes longer.
            \item î, longer yi, \textipa{[i]}: \textit{île}
            \item ô, soft ò, \textipa{[o]}: \textit{hôtel}
            \item û, longer and softer u, \textipa{[y]}: \textit{jeûne}
        \end{enumerate}
    
    \section{Tréma (ë)}
        Vowel has to be pronounced separately from the previous one, same sound as è.
        \begin{itemize}
            \item \textipa{[E]}: \textit{Noël}
        \end{itemize}
    
    \section{Cédille (ç)}
        Usually followed by a, o, u, never an e or i.
        \begin{itemize}
            \item \textipa{[s]}: \textit{garçon}
        \end{itemize}

\chapter{Silent Letters}

    \begin{enumerate}
        \item E preceded by a vowel
            \begin{itemize}
                \item IE = \textipa{[i]}: \textit{pied}
                \item AIE = AI(È), \textit{monnaie}: \textipa{[mOnE]} (change)
                \item UE, \textit{avenue}: \textipa{[av@ny]} 
                \item OIE (wà), \textit{poisson}: \textipa{[pwAsO]}
                \item OUE (wù wà), \textit{boue}: \textipa{[bu]} (bù wà) (mud)
                \item EUE (wù\_è), open as it is the last sound. \textit{vœu}: \textipa{[v2]} (wish)
            \end{itemize}

        \item E preceded by a consonant
            \begin{enumerate}
                \item B, C, D, G, N, P, R, S, T. Do not pronounce the E, it actives the consonant before.
            \end{enumerate}

        \item E in the middle of the word
            \begin{itemize}
                \item No rule, the following are words where e isn't pronounced:
                    \begin{center}
                        un paiement \textipa{[pajm\~{a}]} (payment) \\
                        un remerie \textipa{[r@mEri]} (thanks) \\
                        une proneade \textipa{[pR\~{O}nad]} (walk)
                    \end{center}
                \item E is pronounced in articles and small words (2-3 letters).
            \end{itemize}
        
        \item Last consonant of the word
            \begin{itemize}
                \item B, a few words: \textit{plomb} \textipa{[pl\~{O}]} (lead)
                \item C: \textit{banc} \textipa{[b\~{a}]} (bench)
                \item D: \textit{fond} \textipa{[f\~{O}]} (bottom)
                \item F: un \textit{œuf} \textipa{[2f]} (egg)
                \item G: le poing \textipa{[pw\~{a}]} (fist), except for english words.
                \item L: \textit{outil} \textipa{[uti]} (tool)
                \item P: \textit{coup} \textipa{[ku]} (hit)
                \item S, how words are spelled in sigular: \textit{fils} \textipa{[fis]} (son)
                \item T: \textit{l'art} \textipa{[laR]} (art)
                \item X: only for just how the word is spelled, \textit{paix} \textipa{[pE]} (peace)
                \item Z: \textit{le nez} \textipa{[l@ne]} (nose)
            \end{itemize}

        \item Plural of words
            \begin{itemize}
                \item S: not pronounced: \textit{les chats} \textipa{[leSa]} (cats)
                \item X: not pronounced.
            \end{itemize}
        
        \item Plural of verbs
            \begin{itemize}
                \item -ENT: \textit{il chante, ils chantent} \textipa{[il S\~{a}t]} (sing)
                \item *Irregular vern: \textit{elle dort} \textipa{[El dOr]}, \textit{elles dorment} \textipa{[il dOrm]} (sleep)
            \end{itemize}

        \item The letters H, U, M
            \begin{itemize}
                \item H
                    \begin{enumerate}
                        \item Do not pronounce if it's in the beginning of the word: \textit{hôtel} \textipa{[Otel]} (hotel)
                        \item Acts as separation in the middle of the word: \textit{cohérent} \textipa{[ko.eR\~{a}]} (coherent)
                    \end{enumerate}
                \item U
                    \begin{enumerate}
                        \item When u is after q, do not pronounce u: \textit{qui} \textipa{[ki]} (who)
                        \item When u is after g, do not pronounce u: \textit{guerre} \textipa{[gER]} (war)
                    \end{enumerate}
                \item M
                    \begin{enumerate}
                        \item L'automne \textipa{[lotOm]} (autumn), condamner \textipa{[kO.damne]} (condemn)
                    \end{enumerate}
            \end{itemize}
    \end{enumerate}

\chapter{Links}
    \section{Mandatory Links}
        \begin{enumerate}
            \item Articles and nouns: \textit{les amis} \textipa{[le.zami]} (friends)
            \item Numbers and nouns: \textit{vingt ans} \textipa{[v\~{a}z]} (twenty years)
            \item Adjective and nouns: \textit{Bon anniversaire} \textipa{[bOnanivERsER]} (happy birthday)
            \item Pronouns and pronouns: \textit{Nous y allons} \textipa{[nu.zi.al\~{O}]} (we are going)
            \item Pronouns and verbs: \textit{Nous étuions} \textipa{[nu.zetudj\~{O}]} (we were studying)
            \item After short adverbs: \textit{Bien arrivé} \textipa{[bj\~{a}nari.ve]} (well arrived)
            \item \textit{Quand est-ce que} \textipa{[k\~{a}tEstk@]} (when)
        \end{enumerate}

    \section{Forbidden Links}
        \begin{enumerate}
            \item Before H aspiré (act as consonant), non-silent: \textit{un hariocot} \textipa{[ynarikO]} (bean)
            \item Number before 11
            \item Before \textit{oui}
            \item After a name
            \item After singular nouns
            \item After \textit{et, comment, combien, touhours, quand} with the exception of \textit{quand est-ce que}
        \end{enumerate}

\chapter{Sound}

    \section{é \textipa{[e]}}
        \begin{enumerate}
            \item é, found anywhere in the word or end of passé composé: \textit{été} \textipa{[ete]} (summer)
            \item AI
            \begin{itemize}
                \item future form of verbs with \textit{Je}: \textit{j'irai} \textipa{[ZirE]} (I will go)
                \item some specific words: \textit{aimer} \textipa{[Eme]} (to like)
            \end{itemize} 
            \item ER
                \begin{itemize}
                    \item At the end of nouns: \textit{hiver} \textipa{[ivE]} (winter)
                    \item Infinitive form of -ER verbs: \textit{manger} \textipa{[m\~{a}Ze]} (to eat)
                \end{itemize}
            \item EZ
                \begin{itemize}
                    \item At the end of some words: \textit{assez} \textipa{[asE]} (enough)
                    \item Present form of \textit{vous}: \textit{vous mangez} \textipa{[vu m\~{a}Ze]} (you eat)
                \end{itemize}
            \item ED: \textit{un pied} \textipa{[\~{œ}pje]} (foot)
            \item E
                \begin{itemize}
                    \item At the beginning of words: \textit{Effacer} \textipa{[efase]} (to erase)
                    \item The word \textit{et} by itself is always pronounced \textipa{[e]} (and), different if it's part of a word.
                \end{itemize}
        \end{enumerate}

    \section{è \textipa{[E]}}
        \begin{enumerate}
            \item È, anywhere in the word
            \item Ê, anywhere in the word
            \item Ë: \textit{Noël} \textipa{[nOEl]} (Christmas)
            \item E
                \begin{itemize}
                    \item Before the last audible letters: \textit{la mer} \textipa{[lamE]} (the sea)
                    \item Before two consonants (or more): \textit{celle} \textipa{[sEl]} (that)
                \end{itemize}
            \item AIS
                \begin{itemize}
                    \item At the end of some words: \textit{mais} \textipa{[mE]} (but)
                    \item Imparfait form with \textit{Je. tu}: \textit{je mangeais} \textipa{[Z@ m\~{a}ZE]} (I was eating)
                    \item Conditionel form with \textit{Je, tu}: \textit{je mangerais} \textipa{[Z@ m\~{a}gR@]} (I would eat) \\
                            * If no -s, it means \textit{will}, not \textit{would}
                \end{itemize}
            \item AIT
                \begin{itemize}
                    \item at the end of some words: \textit{du lait} \textipa{[dylE]} (milk)
                    \item Imparfait form with \textit{Il, elle, on}: \textit{il mangeait} \textipa{[il m\~{a}ZE]} (he was eating)
                    \item Conditionel form with \textit{Il, elle, on}: \textit{il mangerait} \textipa{[il m\~{a}gR@]} (he would eat)
                \end{itemize}
            \item AIENT, do not pronounce the \textit{ENT}
                \begin{itemize}
                    \item Imparfait form with \textit{Ils, elles}: \textit{ils mangeaient} \textipa{[il m\~{a}ZE]} (they were eating)
                    \item COnditionel form with \textit{Ils, elles}: \textit{ils mangeraient} \textipa{[il m\~{a}gR@]} (they would eat)
                \end{itemize}
            \item EI, anywhere in the word: \textit{un reine} \textipa{[œR@n]} (queen)
            \item ET, at the end of words: \textit{un livrat} \textipa{[livrE]} (booklet)
            \item AI
                \begin{itemize}
                    \item inside the word: \textit{la laine} \textipa{[lE lEn]} (wool)
                    \item at the end of some words, not related to conjugation: \textit{un balai} \textipa{[bale]} (broom)
                    \item Varies depends on area, might not be light.
                \end{itemize}
            \item ES / EST
                \begin{itemize}
                    \item Present form of \textit{être} with \textit{tu, il, elle, on}: \textit{tu es} \textipa{[ty E]} (you are)
                    \item ES, in small words: \textit{des, les, ces, mes, tes, ses}
                \end{itemize}
            \item EY: \textit{unu poney} \textipa{[yn@ pOnE]} (pony)
            \item ECT, end of words: \textit{un suspect} \textipa{[syspE]} (suspect)
        \end{enumerate}

    \section{Nasal Vowel ON}
        \begin{enumerate}
            \item ON: \textit{Bonbon} \textipa{[b\~{O}b\~{O}]} (candy)
            \item OM, before B or P: \textit{Compter} \textipa{[k\~{O}te]} (to count)
        \end{enumerate}

    \section{AN \textipa{[\~{a}]}}
        \begin{enumerate}
            \item AN: \textit{Dans} \textipa{[d\~{a}]} (in)
            \item AM, before B or P: \textit{Champ} \textipa{[S\~{a}]} (field)
            \item EN: \textit{lent} \textipa{[l\~{a}]} (slow)
            \item EM, before B or P: \textit{le temps} \textipa{[l@ t\~{a}]} (the weather)
            \item * extra: AON, found in three words:
                \begin{enumerate}
                    \item \textit{un faon} \textipa{[f\~{a}]} (fawn)
                    \item \textit{un town} \textipa{[t\~{a}]} (town)
                    \item \textit{un paon} \textipa{[p\~{a}]} (peacock)
                \end{enumerate}
        \end{enumerate}
    
    \section{IN \textipa{[\~{E}]}}
        Higher compare to \textipa{[\~{œ}]}
        \begin{enumerate}
            \item IN: \textit{fin} \textipa{[f\~{E}]} (end)
            \item IM, before B or P: \textit{simple} \textipa{[s\~{E}pl]} (simple)
            \item AIN: \textit{pain} \textipa{[p\~{E}]} (bread)
            \item EIN: \textit{une peinture} \textipa{[yn p\~{E}tyR]} (a painting)
            \item AIM: \textit{la faim} \textipa{[la f\~{E}]} (hunger)
            \item YN: \textit{un lynx} \textipa{[l\~{E}ks]} (lynx)
            \item YM, before B or P: \textit{un symbole} \textipa{[s\~{E}bOl]} (symbol)
            \item EN, irregular but common word: \textit{un examen} \textipa{[Egzam\~{E}]} (exam)
        \end{enumerate}
    
    \section{UN \textipa{\~{œ}}}
        \begin{enumerate}
            \item In northern France, same as \textipa{[\~{E}]} (ā èng)
            \item In southern France, lower sound (èng)
        \end{enumerate}

    \section{IEN \textipa{[j\~{E}]}}
        Rolling I, like third accent: \textit{un bien} \textipa{[j\~{E}]} (good)
    
    \section{OIN \textipa{[w\~{E}]}}
        \textit{Moins} \textipa{[mw\~{E}]} (less)
\end{document}
